The parent paper studies class-uniform rules: row $c$ of a $d\times d$ matrix is a fixed linear function of the matrix
with its columns rotated by $c$, the function depending only on the class of $c$ in a partition of $\Z/d$. Its
Theorem~TP bounds the codes of such a rule by $(d^2+Md)/2$ when the $M$ classes are intervals, and its
Proposition~8.2 shows that the bound can fail for classes with no almost-period: at $d=6$ the classes $\{0,1,3\}$ and
$\{2,4,5\}$ carry a code of dimension $25>24$ over $\F_2$. This memo adds three things.

First, this gain does not come from characteristic~$2$: the same classes carry a code of dimension~$25$ over every
field in which $3\ne0$, in particular over $\Q$, $\R$ and~$\C$.

Second, blowing up a class-uniform code by a factor $k$ multiplies its dimension by $k^2$ and its number of classes by
$k$, so it multiplies the code's excess over the interval value by $k^2$. Applied to the code at $d=6$ it gives, for
$d=6k$ and $2k$ classes, codes whose excess is $k^2$. This answers the parent paper's Problem~9.1(c) as worded, but only
because the number of classes grows: the excess is $dM/12$, of the order of the interval term $dM/2$, and blow-ups
leave the ratio of the excess to $dM$ unchanged. We state the problem for a bounded number of classes.

Third, we price advice in the one-hop model. Advice that is a linear map of rank $q$ costs at most $q$ dimensions,
however each row uses it, and every code extends by exactly $q$ dimensions with suitable advice of rank~$q$. Over a
finite field, advice of any kind with $q$ symbols costs at most $q$ dimensions plus half a row for each class change it
induces. The parent paper suggests charging each bit revealed by a depth reduction as a regime change, at half a row.
In the one-hop model the advice itself costs one dimension per symbol, and only the class changes it induces cost half
a row each.

Every theorem, proposition and corollary is proved in Lean, on top of the parent paper's library; one remark records
computations that the data generator checks. All results concern the linear one-hop model of the parent paper, not
machines.
